
Two numbers characterize the central finding: p = 0.583 and p = 0.0053. The same datasets, the same outcome variable, the same Kaplan-Meier log-rank test — but one analysis uses unadjusted correlations and the other uses propensity-matched pairs. The transformation between these results is attributable to suppressor confounding: high-FAIR datasets tend to be newer, and newer datasets have had less time to accumulate runs, creating a negative confound that masks any genuine FAIR-related effect in unadjusted analysis.

The matched analysis, conducted on a synthetic proof-of-concept cohort (n = 200, 35 matched pairs), yields Cox HR = 3.159 [1.032, 9.672] and a median TTFR difference of 44 days (22\% relative reduction). These results are preliminary. The confidence interval spans nearly an order of magnitude; the cohort is synthetic; the IV is a structural proxy rather than true F-UJI sub-criteria; and a potential PH violation means the average hazard ratio may mischaracterize a time-varying effect. Taken together, these factors make any specific effect size estimate unreliable at this stage.

The methodological contribution is more robust: matched observational designs appear necessary for credible FAIR-outcome studies in ML repositories, because the confounding structure that operates in this data — newer, higher-FAIR datasets having had less time to accumulate runs — would otherwise suppress or distort any genuine FAIR signal. Prior null results in the literature may reflect this confounding rather than the absence of a true effect.

Three experiments remain necessary before the substantive claims in this paper can be considered established: (1) production-scale replication of H-M1 with real OpenML upload dates (obtainable via individual dataset API calls) and true F-UJI sub-criteria scores; (2) execution of H-M3 to test whether the Reusable dimension shows stronger effects on sustained engagement than Findable; and (3) execution of H-M4 to test cross-repository generalization on HuggingFace. The present paper reports the proof-of-concept phase of a larger research program.

